\section{Coherence of the genuine original frame transition units}

\begin{lemma}\label{lem:fc-original-frame-transition-unit-value}
\lean{EnvelopeProperty.Berkovich.SchemeResidueFrameTransition.transitionUnit_val}
\leanok
\uses{lem:fc-original-frame-transition-coefficient,prop:generated-line-generic-frame-norms}
For any commutative ring $A$, original module $P$ on
$\operatorname{Spec}A$ and two actual entire frames $e_0,e_1$, the
value of the genuine transition unit $u(e_0,e_1)$ is the actual
$e_1$-coefficient of the original section $e_0^{-1}(1)$.
Thus the original unit is determined by the actual original frames.
\end{lemma}
\begin{proof}
The original section $e_0^{-1}(1)$ has $e_0$-coefficient one,
by the whole frame inverse identity. Apply the genuine coefficient
comparison to this section. The actual affine section isomorphism
sends one to one, so the resulting identity is
$c_1(e_0^{-1}(1))=u(e_0,e_1)\cdot1$. This proves the assertion.
\end{proof}

\begin{lemma}\label{lem:fc-original-frame-transition-unit-self}
\lean{EnvelopeProperty.Berkovich.SchemeResidueFrameTransition.transitionUnit_self}
\leanok
\uses{lem:fc-original-frame-transition-unit-value}
For every actual entire original frame $e$, one has $u(e,e)=1$.
\end{lemma}
\begin{proof}
The preceding formula expresses the value of this unit as the
coefficient of $e^{-1}(1)$ in $e$, which is one by the actual
whole inverse identity. Equality of values gives equality of units.
\end{proof}

\begin{lemma}\label{lem:fc-original-frame-transition-unit-composition}
\lean{EnvelopeProperty.Berkovich.SchemeResidueFrameTransition.transitionUnit_trans}
\leanok
\uses{lem:fc-original-frame-transition-unit-value,lem:fc-original-frame-transition-coefficient}
For any three actual entire original frames of the same original
module one has
$u(e_0,e_2)=u(e_0,e_1)u(e_1,e_2)$.
\end{lemma}
\begin{proof}
Evaluate the original section $e_0^{-1}(1)$ in $e_2$. The first
lemma identifies this coefficient with $u(e_0,e_2)$. The genuine
coefficient comparison from $e_1$ to $e_2$ gives instead
$u(e_1,e_2)c_1(e_0^{-1}(1))$. The first lemma identifies its remaining
coefficient with $u(e_0,e_1)$. Commutativity of $A$ gives the stated
product in $A$, and equality of values gives equality of units.
\end{proof}

\begin{corollary}\label{cor:fc-original-frame-transition-unit-inverse}
\lean{EnvelopeProperty.Berkovich.SchemeResidueFrameTransition.transitionUnit_symm}
\leanok
\uses{lem:fc-original-frame-transition-unit-self,lem:fc-original-frame-transition-unit-composition}
The reverse actual original frame transition is the inverse unit:
$u(e_1,e_0)=u(e_0,e_1)^{-1}$.
\end{corollary}
\begin{proof}
The composition identity for $e_0,e_1,e_0$ says that the two
transition units multiply to $u(e_0,e_0)=1$. Multiplying by the
inverse of the first unit gives exactly the asserted equality.
\end{proof}
